% GENERATED FILE. Edit canonical lessons, then run npm run generate:books.
% canonical-source-hash: fnv1a64:45b51e2e615e8585
% canonical-lessons: JA-C278-inazuma, JA-C278-ataru, JA-C278-hageshii, JA-C278-tatsumaki, JA-C278-katamuku

\chapter{Lightning}
\label{ch:ja-lightning}
\hlchaptermodality{278}

\begin{chapteropening}
\textbf{By the end of this chapter:} \emph{I can say what a storm struck and blew down.}

Complete the last lesson of chapter 278: I can say what a storm struck and blew down.
\end{chapteropening}

\section[\texorpdfstring{\ja{いなずま} (inazuma)}{いなずま (inazuma)}]{\ja{いなずま} (inazuma) — a flash of lightning}

\label{lesson:JA-C278-inazuma}



\begin{quote}
Before the new one: say the Japanese for a power cut, then the Japanese for to go numb with cold.
\end{quote}

\subsection*{You'll want to know: \ja{いなずま}}
\textbf{\ja{いなずま}} — \emph{inazuma} — \textquotedblleft{}a flash of lightning\textquotedblright{}.

\begin{quote}
\textbf{\ja{ゆうべ、そらに} \ja{いなずまが} \ja{ひかりました。}} — \emph{yūbe, sora ni inazuma ga hikarimashita} — \textquotedblleft{}Last night lightning flashed across the sky.\textquotedblright{}
\end{quote}

\begin{cousinweb}
\textbf{\ja{いなずま}} (\emph{inazuma}) is a native Japanese word (wago), inherited rather than borrowed from Chinese; its kanji was attached later as a kun'yomi (JMdict).
\end{cousinweb}

\begin{grammarlens}[title={What you've built}]
One more word for telling what happened.
\end{grammarlens}

\subsection*{Your turn}
\begin{itemize}
\raggedright
  \item \emph{Say it:} \emph{inazuma}
  \item \emph{Say it:} \emph{inazuma}, once more
  \item \emph{Say it:} \emph{yūbe, sora ni inazuma ga hikarimashita} — Last night lightning flashed across the sky.
  \item \emph{Recall:} say the Japanese for a power cut, then the Japanese for to go numb with cold, then say \emph{inazuma} again
\end{itemize}

\begin{tcolorbox}[breakable,colback=teal!4,colframe=teal!35!black,title={Before you move on}]
What does \ja{いなずま} mean? (\textquotedblleft{}A flash of lightning\textquotedblright{}.) Say it once more.
\end{tcolorbox}
\section[\texorpdfstring{\ja{あたる} (ataru)}{あたる (ataru)}]{\ja{あたる} (ataru) — to strike; to hit (a target)}

\label{lesson:JA-C278-ataru}



\begin{quote}
Before the new one: say the Japanese for to go numb with cold, then the Japanese for a flash of lightning.
\end{quote}

\subsection*{You'll want to know: \ja{あたる}}
\textbf{\ja{あたる}} — \emph{ataru} — \textquotedblleft{}to strike; to hit (a target)\textquotedblright{}.

\begin{quote}
\textbf{\ja{いしが} \ja{くるまの} \ja{まどに} \ja{あたりました。}} — \emph{ishi ga kuruma no mado ni atarimashita} — \textquotedblleft{}A stone hit the car window.\textquotedblright{}
\end{quote}

Polite \textbf{\ja{あたります}}, \emph{atarimasu}; a godan verb. Past (polite): \textbf{\ja{あめが} \ja{まどに} \ja{あたりました。}}, \emph{ame ga mado ni atarimashita}.

\begin{cousinweb}
\textbf{\ja{あたる}} (\emph{ataru}) is a native Japanese word (wago), inherited rather than borrowed from Chinese; its kanji was attached later as a kun'yomi (JMdict).
\end{cousinweb}

\begin{grammarlens}[title={What you've built}]
One more word for telling what happened.
\end{grammarlens}

\subsection*{Your turn}
\begin{itemize}
\raggedright
  \item \emph{Say it:} \emph{ataru}
  \item \emph{Say it:} \emph{ataru}, once more
  \item \emph{Say it:} \emph{ishi ga kuruma no mado ni atarimashita} — A stone hit the car window.
  \item \emph{Recall:} say the Japanese for to go numb with cold, then the Japanese for a flash of lightning, then say \emph{ataru} again
\end{itemize}

\begin{tcolorbox}[breakable,colback=teal!4,colframe=teal!35!black,title={Before you move on}]
What does \ja{あたる} mean? (\textquotedblleft{}To strike; to hit (a target)\textquotedblright{}.) Say it once more.
\end{tcolorbox}
\section[\texorpdfstring{\ja{はげしい} (hageshii)}{はげしい (hageshii)}]{\ja{はげしい} (hageshii) — violent; fierce}

\label{lesson:JA-C278-hageshii}



\begin{quote}
Before the new one: say the Japanese for a flash of lightning, then the Japanese for to strike.
\end{quote}

\subsection*{You'll want to know: \ja{はげしい}}
\textbf{\ja{はげしい}} — \emph{hageshii} — \textquotedblleft{}violent; fierce\textquotedblright{}.

\begin{quote}
\textbf{\ja{きのう} \ja{日本中}\ja{で} \ja{はげしい} \ja{あめが} \ja{ふりました。}} — \emph{kinō nihonjū de hageshii ame ga furimashita} — \textquotedblleft{}Yesterday fierce rain fell all over Japan.\textquotedblright{}
\end{quote}

An i-adjective.

\begin{cousinweb}
\textbf{\ja{はげしい}} (\emph{hageshii}) is a native Japanese word (wago), inherited rather than borrowed from Chinese; its kanji was attached later as a kun'yomi (JMdict).
\end{cousinweb}

\begin{grammarlens}[title={What you've built}]
One more word for telling what happened.
\end{grammarlens}

\subsection*{Your turn}
\begin{itemize}
\raggedright
  \item \emph{Say it:} \emph{hageshii}
  \item \emph{Say it:} \emph{hageshii}, once more
  \item \emph{Say it:} \emph{kinō nihonjū de hageshii ame ga furimashita} — Yesterday fierce rain fell all over Japan.
  \item \emph{Recall:} say the Japanese for a flash of lightning, then the Japanese for to strike, then say \emph{hageshii} again
\end{itemize}

\begin{tcolorbox}[breakable,colback=teal!4,colframe=teal!35!black,title={Before you move on}]
What does \ja{はげしい} mean? (\textquotedblleft{}Violent; fierce\textquotedblright{}.) Say it once more.
\end{tcolorbox}
\section[\texorpdfstring{\ja{たつまき} (tatsumaki)}{たつまき (tatsumaki)}]{\ja{たつまき} (tatsumaki) — a tornado; a whirlwind}

\label{lesson:JA-C278-tatsumaki}



\begin{quote}
Before the new one: say the Japanese for to strike, then the Japanese for violent.
\end{quote}

\subsection*{You'll want to know: \ja{たつまき}}
\textbf{\ja{たつまき}} — \emph{tatsumaki} — \textquotedblleft{}a tornado; a whirlwind\textquotedblright{}.

\begin{quote}
\textbf{\ja{たつまきで} \ja{やねが} \ja{とびました。}} — \emph{tatsumaki de yane ga tobimashita} — \textquotedblleft{}The tornado blew the roof off.\textquotedblright{}
\end{quote}

\begin{cousinweb}
\textbf{\ja{たつまき}} (\emph{tatsumaki}) is a native Japanese word (wago), inherited rather than borrowed from Chinese, built from \emph{tatsu}, \textquotedblleft{}dragon\textquotedblright{}, and \emph{maki}, \textquotedblleft{}winding\textquotedblright{} (from \emph{maku}) (JMdict, KANJIDIC2).
\end{cousinweb}

\begin{grammarlens}[title={What you've built}]
One more word for telling what happened.
\end{grammarlens}

\subsection*{Your turn}
\begin{itemize}
\raggedright
  \item \emph{Say it:} \emph{tatsumaki}
  \item \emph{Say it:} \emph{tatsumaki}, once more
  \item \emph{Say it:} \emph{tatsumaki de yane ga tobimashita} — The tornado blew the roof off.
  \item \emph{Recall:} say the Japanese for to strike, then the Japanese for violent, then say \emph{tatsumaki} again
\end{itemize}

\begin{tcolorbox}[breakable,colback=teal!4,colframe=teal!35!black,title={Before you move on}]
What does \ja{たつまき} mean? (\textquotedblleft{}A tornado; a whirlwind\textquotedblright{}.) Say it once more.
\end{tcolorbox}
\section[\texorpdfstring{\ja{かたむく} (katamuku)}{かたむく (katamuku)}]{\ja{かたむく} (katamuku) — to lean; to tilt}

\label{lesson:JA-C278-katamuku}



\begin{quote}
Before the new one: say the Japanese for violent, then the Japanese for a tornado.
\end{quote}

\subsection*{You'll want to know: \ja{かたむく}}
\textbf{\ja{かたむく}} — \emph{katamuku} — \textquotedblleft{}to lean; to tilt\textquotedblright{}.

\begin{quote}
\textbf{\ja{じしんで} \ja{たなが} \ja{かたむいています。}} — \emph{jishin de tana ga katamuite imasu} — \textquotedblleft{}The shelves are leaning since the earthquake.\textquotedblright{}
\end{quote}

Polite \textbf{\ja{かたむきます}}, \emph{katamukimasu}; a godan verb. Past (polite): \textbf{\ja{きが} \ja{かたむきました。}}, \emph{ki ga katamukimashita}.

\begin{cousinweb}
\textbf{\ja{かたむく}} (\emph{katamuku}) is a native Japanese word (wago), inherited rather than borrowed from Chinese; its kanji was attached later as a kun'yomi (JMdict).
\end{cousinweb}

\begin{grammarlens}[title={What you've built}]
One more word for telling what happened.
\end{grammarlens}

\subsection*{Your turn}
\begin{itemize}
\raggedright
  \item \emph{Say it:} \emph{katamuku}
  \item \emph{Say it:} \emph{katamuku}, once more
  \item \emph{Say it:} \emph{jishin de tana ga katamuite imasu} — The shelves are leaning since the earthquake.
  \item \emph{Recall:} say the Japanese for violent, then the Japanese for a tornado, then say \emph{katamuku} again
\end{itemize}

\begin{tcolorbox}[breakable,colback=teal!4,colframe=teal!35!black,title={Before you move on}]
What does \ja{かたむく} mean? (\textquotedblleft{}To lean; to tilt\textquotedblright{}.) Say it once more.
\end{tcolorbox}
